\documentclass[12pt, a4paper]{article}

\usepackage[utf8]{inputenc}
\usepackage[T1]{fontenc}
\usepackage[brazil]{babel}
\usepackage[top=2.5cm, bottom=2.5cm, left=3cm, right=3cm]{geometry}
\usepackage{booktabs}
\usepackage{array}
\usepackage{parskip}
\usepackage{amssymb}

\begin{document}

\begin{center}
    \textbf{DEPARTAMENTO DE INFORMÁTICA E ESTATÍSTICA - CTC - UFSC}\\[0.4em]
    \textbf{RATIFICAÇÃO DE PLANO DE TRABALHO DO SEMESTRE PARA DESENVOLVIMENTO DE TCC}
\end{center}

\vspace{1em}

\noindent\textbf{Disciplina:} \quad ($\square$) TCC1 \quad ($\boxtimes$) TCC2

\noindent\textbf{Curso:} \quad ($\boxtimes$) CCO \quad ($\square$) SIN

\vspace{1em}

\noindent\textbf{Nome do(s) autor(es):} Artur Luiz Rizzato Toru Soda

\vspace{1em}

\noindent\textbf{Título do TCC:} Estudo de Viabilidade do uso de Instâncias Burstable e Spot da AWS para Computação de Alto Desempenho

\vspace{1.5em}

\noindent\textbf{Objetivos:}

\noindent Concluir o TCC integrando as duas contribuições desenvolvidas durante o período de IC:
(1) análise de viabilidade de instâncias burstable, publicada no WCC 2025; e
(2) implementação e avaliação de tolerância a falhas transparente em instâncias spot com MANA.
Como extensão, implementar uma política adaptativa de recuperação no HPC@Cloud, uma
terceira estratégia (ADAPTIVE) que seleciona automaticamente entre REPLACE e DEGRADED
com base no tamanho do cluster e no tempo decorrido da execução, derivando a heurística
de decisão da análise de crossover já realizada, e validá-la experimentalmente.

\vspace{1.5em}

\noindent\textbf{Cronograma (Atividades do Semestre):}

\vspace{0.5em}

\noindent
\begin{tabular}{@{} p{3cm} p{10cm} @{}}
    \toprule
    \textbf{Período} & \textbf{Atividades} \\
    \midrule
    Agosto    & Retificação do plano de trabalho; Relatório de IC \\
    Setembro  & Implementação da política adaptativa; Experimentos de validação \\
    Outubro   & Escrita da monografia final (integração das duas partes e política adaptativa) \\
    Novembro  & Revisões e ajustes finais \\
    Dezembro  & Defesa do TCC \\
    \bottomrule
\end{tabular}

\vspace{2em}

\noindent\textbf{Nome do Professor responsável:} Prof. Márcio Bastos Castro, Dr.

\end{document}
